\Section{Определения предела числовой функции одного переменного в
терминах окрестностей и в терминах последовательностей, их эквивалентность.
Свойства пределов функции. Критерий Коши существования конечного предела
функции. Предел сложной функции. Существование
односторонних пределов у монотонных функций.}

\Subsection{Определения предела числовой функции одного переменного
в терминах окрестностей и в терминах последовательностей, их эквивалентность}

\begin{definition}
Пусть даны точка \(x_{0}\) и действительное число \(\delta > 0\). Тогда интервал
\((x_{0} - \delta, ~ x_{0} + \delta)\) называют \textit{\(\delta\)-окрестностью точки \(x_{0}\).}
\[
x \in U_{\delta}(x_{0}) ~ \Leftrightarrow ~ x \in (x_{0} - \delta, ~ x_{0} + \delta)
~ \Leftrightarrow ~ |x - x_{0}| < \delta
\]
\end{definition}

\begin{definition}
Пусть даны точка \(x_{0}\) и действительное число \(\delta > 0\).
Тогда объединение интервалов \((x_{0} - \delta, ~ x_{0}) \cup (x_{0}, ~ x_{0} + \delta)\)
называют \textit{проколотой \(\delta\)-окрестностью точки \(x_{0}\).}
\[
x \in \overset{\circ}{U}_{\delta}(x_{0}) ~ \Leftrightarrow ~
x \in (x_{0} - \delta, ~ x_{0}) \cup (x_{0}, ~ x_{0} + \delta)
~ \Leftrightarrow ~ 0 < |x - x_{0}| < \delta
\]
\end{definition}

Далее мы считаем, что функция \(f(x)\) определена в
некоторой проколотой \(\Delta\)-окрестности точки \(x_{0}\).

\begin{definition}[По Коши] \label{def_lim_Koshi}
Число \(\alpha\) называют \textit{пределом функции \(f(x)\) в точке \(x_{0}\)},
если для любого числа \(\epsilon > 0\) найдется такое число \(\delta > 0\),
что для всех \(x\) из проколотой \(\delta\)-окрестности точки \(x_{0}\)
соответствующие им значения \(f(x)\) принадлежат \(\epsilon\)-окрестности точки \(\alpha\).
\[
[\lim_{ x \to x_{0} } f(x) = \alpha ] \overset{def}{=} \Big[ \forall \, \epsilon > 0
~~ \exists \, \delta > 0 ~~ \forall \, x \in \overset{\circ}{U}_\delta(x_{0})
~ \longmapsto ~ f(x) \in U_{\epsilon}(\alpha) \Big]
\]
\end{definition}

\begin{definition}[По Гейне]
Число \(\alpha\) называют \textit{пределом функции \(f(x)\) в точке \(x_{0}\)},
если для любой последовательности \(\{x_{n}\}\), сходящейся к \(x_{0}\), такой,
что \(x_{n} \neq x_{0}\) для всех \(n \in \N\), последовательность
\(\{f(x_{n})\}\) сходится к числу \(\alpha\).
\[
[\lim_{ x \to x_{0} } f(x) = \alpha ] \overset{def}{=} \Big[ \forall \, \{x_{n}\}:
[\lim_{ n \to \infty } x_{n} = x_{0}] ~ \& ~ [\forall \, n \in \N ~~ x_{n} \neq x_{0}]
~ \longmapsto ~ \lim_{ n \to \infty } f(x_{n}) = \alpha \Big]
\]
\end{definition}

\begin{note}
Будем называть последовательности, которые удовлетворяют условиям:
\[
\lim\limits_{n \to \infty}{x_{n}} = x_{0} \qquad \&
\qquad \forall \, n \in \N ~~ x_{n} \neq x_{0}
\]
\textit{последовательностями Гейне}.
\end{note}

\begin{definition}
Аналогично пределу функции в точке определяются \textit{пределы функции в точке слева и справа}
(также называемые \textit{односторонними пределами}):
\[
[f(x_{0} - 0) = \alpha] \overset{def}{=} \Big[ \forall \, \epsilon > 0 ~~ \exists \, \delta > 0 ~~
\forall \, x: ~~ 0 < x_{0} - x < \delta ~ \longmapsto ~ |f(x) - \alpha| < \epsilon \Big]
\]
\[
[f(x_{0} + 0) = \alpha] \overset{def}{=} \Big[ \forall \, \epsilon > 0 ~~ \exists \, \delta > 0 ~~
\forall \, x: ~~ 0 < x - x_{0} < \delta ~ \longmapsto ~ |f(x) - \alpha| < \epsilon \Big]
\]
\end{definition}

\begin{statement} \label{statement_side_lims}
Если в точке \(x_{0}\) левый и правый пределы функции \(f(x)\) совпадают,
то в точке \(x_{0}\) существует предел \(f(x)\), равный односторонним пределам.

\medskip

\proofname~Пусть односторонние пределы равны \(\alpha\). Тогда, согласно только что
обозначенным определениям, для произвольного \(\epsilon > 0\) должно выполняться следующее:
\[
\exists \, \delta_{1} > 0 ~~\forall \, x: ~~ 0 < x_{0} - x < \delta_{1}
~ \longmapsto ~ |f(x) - \alpha| < \epsilon,
\]
\[
\exists \, \delta_{2} > 0 ~~\forall \, x: ~~ 0 < x - x_{0} < \delta_{2}
~ \longmapsto ~ |f(x) - \alpha| < \epsilon.
\]
Выберем \(\delta = \min(\delta_{1}, \, \delta_{2})\). Если переписать вышенаписанное
через \(\delta\), то т.к. оба выражения будут справедливы, получим:
\[
\exists \, \delta > 0 ~~\forall \, x: ~~ 0 < |x_{0} - x| < \delta
~ \longmapsto ~ |f(x) - \alpha| < \epsilon,
\]
то есть по определению число \(\alpha\)~--- предел функции \(f(x)\) в точке \(x_{0}\).~\qedsymbol
\end{statement}

\begin{theorem}
Определения предела функции по Коши и по Гейне эквивалентны.

\medskip

\proofname~Сначала докажем, что предел по Коши является пределом по Гейне. \\
Пусть \(\alpha\) является пределом по Коши функции \(f(x)\) в точке \(x_{0}\).
Пусть \(\{x_{n}\}\)~--- произвольная последовательность, сходящаяся к \(x_{0}\)
и каждый член которой не равен \(x_{0}\). \\

Произвольному числу \(\epsilon > 0\) соответствует некоторое число \(\delta > 0\)
\text{(по Определению \ref{def_lim_Koshi})}, а уже числу \(\delta\) по определению предела \textit{последовательности}
соответствует некоторый номер \(N\), начиная с которого для всех номеров \(n \geq N\)
выполняется \(|x_{n} - x_{0}| < \delta\). Тогда (с учетом условия \(x_{n} \neq x_{0}\))
получаем, что для тех же номеров \(x_{n} \in \overset{\circ}{U}_{\delta}(x_{0})\).
Тогда по определению Коши для них же верно и \(|f(x_{n}) - \alpha| < \epsilon\),
то есть \(\lim\limits_{n \to \infty}{f(x_{n})} = \alpha\).

Поскольку мы брали произвольную последовательность Гейне \(\{x_{n}\}\), то доказанное
верно для любой последовательности Гейне. Следовательно, число \(\alpha\)
является пределом по Гейне.\\
\\
Теперь покажем, что предел по Гейне является пределом по Коши.\\
Предположим противное: пусть число \(\alpha\) является пределом по Гейне,
но не является пределом по Коши:
\[
\forall \{x_{n}\} : ~~ \lim_{ n \to \infty } x_{n} = x_{0} \hspace{6mm} \& \hspace{6mm}
\forall \, n \in \N ~~ x_{n} \neq x_{0} ~ \longmapsto ~ \lim_{ n \to \infty }f(x_{n}) = \alpha
\]
\[ \exists \, \epsilon_{0} > 0 ~~ \forall \, \delta > 0 ~~ \exists \, x_{\delta}
\in \overset{\circ}{U}_{\delta} (x_{0}) : ~~ |f(x_{\delta}) - \alpha| \geq \epsilon_{0}
\]

В таком случае (по второй строке) для положительных чисел \(\delta(n) = \frac{1}{n}\)
будет находиться число \(x_{n}\) такое, что:
\[
0 < |x_{n} - x_{0}| < \frac{1}{n} \qquad \& \qquad|f(x_{n}) - \alpha| \geq \epsilon_{0}
\]

Заметим, что левая часть означает, что \(\lim\limits_{n \to \infty}{x_{n}} = x_{0}\)
и при этом \(x_{n} \neq x_{0}\) при любых натуральных \(n\).
Следовательно, \(\{x_{n}\}\) является последовательностью Гейне, и тогда последовательность
\(\{f(x_{n})\}\) должна сходиться к \(\alpha\). Но это противоречит правой стороне
записанного выше утверждения. Тогда получаем, что предположение неверно, и тогда
предел по Гейне должен являться пределом по Коши.~\qedsymbol
\end{theorem}

\Subsection{Свойства пределов функций}

\begin{theorem} \label{theorem_lim_arithmetic}
Пусть функции \(f(x)\) и \(g(x)\) таковы, что \(\lim\limits_{x \to x_{0}}f(x) = \alpha\)
и \(\lim\limits_{x \to x_{0}} g(x) = \beta\). Тогда функции \(y = f(x) \pm g(x)\),
\(y=f(x)\cdot g(x)\) и \(y=\dfrac{f(x)}{g(x)}\) имеют пределы в точке \(x_{0}\),
равные соответственно \(\alpha \pm \beta\), \(\alpha \cdot \beta\) и \(\dfrac{\alpha}{\beta}\)
(последнее~--- при \(\beta \neq 0\)).

\medskip

\textit{Примечание: эти и следующие свойства доказываются одинаково: мы используем определение
по Гейне, т.е. переходим от функций к последовательностям, используем аналогичные свойства
последовательностей (см. подтему \ref{subsection_lim_arithmetic}) и возвращаемся обратно к функциям. Здесь будет приведено
только доказательство для суммы (разности) функций как образца доказательства остальных свойств.} \\

\proofname~Пусть \(\{x_{n}\}\) - произвольная последовательность Гейне, сходящаяся к \(x_{0}\).
По определению предела \textit{функции} по Гейне
\(\lim\limits_{n \to \infty}{f(x_{n})} = \alpha\) и
\(\lim\limits_{n \to \infty}{g(x_{n})} = \beta\). Тогда для функции \(y = f(x) \pm g(x)\)
по свойствам пределов \textit{последовательностей} имеем:
\[
\lim_{ n \to \infty }(f(x_{n}) \pm g(x_{n})) = \lim_{ n \to \infty } f(x_{n}) \pm
\lim_{ n \to \infty } g(x_{n}) = \alpha \pm \beta.
\]

Поскольку мы брали произвольную последовательность Гейне \(\{x_{n}\}\), то же самое
будет верно для любой последовательности Гейне, сходящейся к \(x_{0}\).
Тогда по определению предела \textit{функции} по Гейне
\(\lim\limits_{x \to x_{0}}{(f(x) \pm g(x))} = \alpha \pm \beta\),
что и требовалось доказать.~\qedsymbol
\end{theorem}

\begin{theorem}
Пусть функция \(f(x)\) определена в некоторой проколотой \(\delta\)-окрестности
точки \(x_{0}\), причем \(\lim\limits_{x \to x_{0}} f(x) = \alpha\) и в той же окрестности
выполнено неравенство \({f(x) \geq \gamma ~ \Big[ f(x) \leq \gamma \Big]}\). Тогда справедливо
неравенство \(\alpha \geq \gamma ~ \Big[ \alpha \leq \gamma \Big]\).
\end{theorem}

\begin{note}
Если заменить нестрогие неравенства \((f(x) \geq \gamma)\) на строгие \((f(x) > \gamma)\),
то для пределов неравенство все равно останется нестрогим \((\alpha \geq \gamma)\).
\end{note}

\begin{corollary}
Пусть функции \(f(x)\) и \(g(x)\) определены в некоторой проколотой
\text{\(\delta\)-окрестности} точки \(x_{0}\), причем
\(\lim\limits_{x\to x_{0}} f(x) = \alpha\),
\(\lim\limits_{x\to x_{0}} g(x) = \beta\) и в той же окрестности выполнено неравенство
\(f(x) \geq g(x)\). Тогда справедливо неравенство \(\alpha \geq \beta\).
\end{corollary}

\begin{theorem}
Пусть в некоторой проколотой \(\delta\)-окрестности определены функции \(f(x)\), \(g(x)\)
и \(h(x)\), причем в этой окрестности выполнены неравенства \(f(x) \leq g(x) \leq h(x)\),
и при этом \(\lim\limits_{x \to x_{0}} f(x) = \lim\limits_{x \to x_{0}} g(x)\).
Тогда справедливо следующее:
\[
\lim_{ x \to x_{0} }f(x) = \lim_{ x \to x_{0} } g(x) = \lim_{ x \to x_{0} }h(x).
\]
\end{theorem}

\begin{note}
Да, по сути это теорема о трех последовательностях, только для функций.
Причем если заменить \(x_{0}\) на \(\infty, ~ \pm \infty\) или \(x_{0} \pm 0\)
(пределы слева и справа), то теорема останется справедливой.
\end{note}

\Subsection{Предел сложной функции}

\begin{theorem}[Предел суперпозиции функций]
Пусть функция \(f(x)\) определена в некоторой проколотой окрестности точки \(x_{0}\),
имеет предел \(\lim\limits_{x \to x_{0}} f(x) = \alpha\) и для всех \(x\) из указанной окрестности
\(f(x) \neq \alpha\). Также предположим, что функция \(g(y)\) определена в некоторой проколотой
окрестности точки \(\alpha\) и \(\lim\limits_{y \to \alpha} g(y) = \beta\).
Тогда справедливо равенство \(\lim\limits_{x \to x_{0}} g(f(x)) = \beta\).

\medskip

\proofname~Пусть \(\{x_{n}\}\) - произвольная последовательность Гейне, сходящаяся к \(x_{0}\).
Тогда, согласно определению по Гейне, \(\lim\limits_{n \to \infty} f(x_{n}) = \alpha\).
Поскольку по условию \(f(x) \neq \alpha\), то последовательность значений функции
\(\{f(x_{n})\}\) сама является последовательностью Гейне, сходящейся к \(\alpha\).
Тогда, согласно определению по Гейне, \(\lim\limits_{n \to \infty} g(f(x_{n})) = \beta\).
Поскольку мы брали произвольную последовательность Гейне \(\{x_{n}\}\), сходящуюся к \(x_{0}\),
то же самое будет верно для любой последовательности Гейне, сходящейся к \(x_{0}\).
Тогда по Гейне \(\lim\limits_{x \to x_{0}} g(f(x)) = \beta\).~\qedsymbol
\end{theorem}

\Subsection{Критерий Коши существования конечного предела функции}

\begin{definition}
Будем говорить, что функция \(f(x)\) в точке \(x_{0}\) удовлетворяет условию Коши,
если для любого \(\epsilon > 0\) найдется такое число \(\delta > 0\), что для всех точек
\(x'\) и \(x''\), принадлежащих проколотой \(\delta\)-окрестности точки \(x_{0}\),
выполняется неравенство \(|f(x') - f(x'')| < \epsilon\).
\[
\begin{aligned}
\forall \, \epsilon > 0 ~~ \exists \, \delta > 0 : ~~ \forall x', ~ x'':
~~ 0 < |x' - x_{0}| < \delta,\\
0 < |x'' - x_{0}| < \delta ~ \longmapsto& ~ |f(x') - f(x'')| < \epsilon.
\end{aligned}
\]
\end{definition}

\begin{theorem}
Для того чтобы функция \(f(x)\) в точке \(x_{0}\) имела \textit{конечный} предел,
необходимо и достаточно, чтобы в этой точке она удовлетворяла условию Коши.

\medskip

\proofname~Докажем необходимость. \\
Пусть функция \(f(x)\) в точке \(x_{0}\) имеет предел:
\(\lim\limits_{x\to x_{0}} f(x) = \alpha\). Тогда для любого \(\epsilon > 0\) найдется
такое \(\delta > 0\), что для любых точек \(x'\) и \(x''\) из проколотой \(\delta\)-окрестности
точки \(x_{0}\) выполнено:
\[
|f(x') - \alpha| < \frac{\epsilon}{2} \qquad |f(x'') - \alpha| < \frac{\epsilon}{2}.
\]

Тогда справедливо следующее:
\[
|f(x') - f(x'')| = |(f(x') - \alpha) - (f(x'') - \alpha)| \leq
|f(x') - \alpha| + |f(x'') - \alpha| < \frac{\epsilon}{2} + \frac{\epsilon}{2} = \epsilon.
\]
Следовательно, по определению функция \(f(x)\) в точке \(x_{0}\) удовлетворяет условию Коши.\\
\\
Теперь докажем достаточность.\\
Пусть функция \(f(x)\) в точке \(x_{0}\) удовлетворяет условию Коши.
Возьмем произвольную последовательность Гейне \(\{x_{n}\}\), сходящуюся к \(x_{0}\).
Докажем, что последовательность \(\{f(x_{n})\}\) сходится.\\
Возьмем произвольное \(\epsilon > 0\) и соответствующее ему по условию Коши \(\delta > 0\).
Заметим, что так как \(\lim\limits_{ n \to \infty } x_{n} = x_{0}\), то, начиная с некоторого номера
\(N\) (\(n \geq N\)), выполнено неравенство \(0 < |x_{n} - x_{0}| < \delta\).
При этом очевидно, что тогда для любого \(p \in \N\) выполнено
\(0 < |x_{n+p} - x_{0}| < \delta\). По условию Коши имеем:
\[
|f(x_{n+p}) - f(x_{n})| < \epsilon.
\]
Следовательно, последовательность \(\{f(x_{n})\}\) фундаментальна.
Тогда по Критерию Коши \textit{для последовательностей} \(\{f(x_{n})\}\) сходится.\\

Осталось доказать, что такие последовательности функций сходятся к одному числу,
а не разным.
Пусть \(\{x'_{n}\}\) и \(\{x''_{n}\}\) - последовательности Гейне, сходящиеся
к \(x_{0}\). Согласно уже доказанному, соответствующие им последовательности функций
должны сходиться. Обозначим:
\[
\lim_{ n \to \infty } f(x'_{n}) = \alpha' \qquad \lim_{ n \to \infty } f(x''_{n}) = \alpha''.
\]

Введем новую последовательность:
\[
\{\overset{\_}{x}_{n}\} = \{x'_{1}, \, x''_{1}, \, x'_{2}, \, x''_{2}, \, \dots\}
\]
Заметим, что \(\{\overset{\_}{x}_{n}\}\) также является последовательностью Гейне,
сходящейся к \(x_{0}\).
Тогда (опять же по вышедоказанному) последовательность \(\{f(\overset{\_}{x}_{n})\}\) сходится.
Обозначим \(\lim\limits_{n \to \infty} f(\overset{\_}{x}_{n}) = \alpha\).
Тогда (по Теореме \ref{theorem_subseries_1}):
\[
\alpha' = \alpha'' = \alpha.
\]

Следовательно, мы доказали, что \textit{все} последовательности \(f(x_{n})\) сходятся
к одному и тому же числу. Тогда по Гейне функция \(f(x)\) в точке \(x_{0}\) имеет
конечный предел, что и требовалось доказать.~\qedsymbol
\end{theorem}

\Subsection{Существование односторонних пределов у монотонных функций}

\begin{theorem}[Теорема о существовании односторонних пределов у монотонной функции]
\label{theorem_sideLims_monotonous}
Если функция \(f(x)\) определена и монотонна на интервале \((a, \, b)\), то в каждой точке
\(x_{0}\) этого интервала она имеет односторонние пределы. Причем если функция
является неубывающей на интервале \((a, \, b)\), то справедливы неравенства
\[
f(x_{0} - 0) \leq f(x_{0}) \leq f(x_{0} + 0),
\]
а если \(f(x)\)~--- невозрастающая на интервале, то справедливы неравенства
\[
f(x_{0} - 0) \geq f(x_{0}) \geq f(x_{0} + 0).
\]

\medskip

\proofname~Пусть \(f(x)\) неубывающая на интервале \((a, \, b)\)
(случай невозрастающей функции аналогичен). \\
Зафиксируем точку \(x_{0} \in (a, \, b)\). Поскольку функция неубывающая,
то для любой точки \(x \in (a, \, x_{0})\) выполняется неравенство \(f(x) \leq f(x_{0})\).
Следовательно, функция \(f(x)\) ограничена сверху на интервале \((a, \, x_{0})\), и тогда
по Теореме \ref{theorem_set_SupInf} существует точная верхняя грань функции \(f(x)\)
на этом интервале:
\[
\alpha = \underset{(a, \, x_{0})}{\sup} f(x).
\]

По определению точной верхней грани для любого \(\epsilon > 0\) найдется такое число
\(x_{\epsilon} \in (a, \, x_{0})\), что справедливы неравенства
\[
\alpha - \epsilon < f(x_{\epsilon}) \leq \alpha.
\]
Введем обозначение \(\delta = x_{0} - x_{\epsilon}.\) Заметим следующее:
\[
x_{\epsilon} \in (a, \, x_{0}) ~ \Rightarrow ~ x_{\epsilon} < x_{0} ~ \Rightarrow ~ \delta > 0,
\]
\[
f(x) - \text{неубывающая} ~ \Rightarrow ~ \alpha - \epsilon < f(x_{\epsilon}) \leq f(x) \leq \alpha
~ \Rightarrow ~ |f(x) - \alpha| < \epsilon.
\]
Собирая все вместе, получаем следующее:
\[
\forall \, \epsilon > 0 ~~ \exists \, \delta > 0 ~~ \forall x: ~ 0 < x_{0} - x < \delta
~ \longmapsto ~ |f(x) - \alpha| < \epsilon.
\]

Тогда по определению предела \(f(x_{0} - 0) = \alpha \leq f(x_{0})\).
Существование предела справа и соответствующее неравенство доказываются аналогично.~\qedsymbol
\end{theorem}

\begin{note}
Теорема остается справедливой и для бесконечных интервалов~--- например, \((a, \, +\infty)\)
или на всей \(\R\).
\end{note}